\documentclass[11pt,letterpaper]{article}
\usepackage[margin=0.75in]{geometry}
\usepackage{fontspec}
\setmainfont{TeX Gyre Termes}
\usepackage{xcolor}
\definecolor{linkblue}{HTML}{244A70}
\usepackage[colorlinks=true,urlcolor=linkblue,linkcolor=black]{hyperref}
\hypersetup{pdftitle={Ruoshui Li - Curriculum Vitae},pdfauthor={Ruoshui Li}}
\usepackage{enumitem}
\usepackage{titlesec}
\usepackage{microtype}
\pagestyle{empty}
\setlength{\parindent}{0pt}
\setlength{\parskip}{2pt}
\setlist[itemize]{leftmargin=14pt,itemsep=1pt,parsep=0pt,topsep=2pt}
\titleformat{\section}{\large\bfseries}{}{0pt}{}
\titlespacing*{\section}{0pt}{9pt}{4pt}
\newcommand{\entry}[3]{\par\textbf{#1}\hfill #2\\#3\par}
\begin{document}
{\LARGE\bfseries Ruoshui (Walter) Li}\par
Department of Mathematics and Statistics, Saint Louis University\par
\href{mailto:ruoshuiliwalter.li@slu.edu}{ruoshuiliwalter.li@slu.edu}
\quad $\vert$ \quad
\href{mailto:liruoshui0710@gmail.com}{liruoshui0710@gmail.com}\par
\textbf{Academic portfolio:} \href{https://github.com/Walter-Li17/academic-projects}{github.com/Walter-Li17/academic-projects}

\section*{Education}
\entry{Saint Louis University}{Fall 2025 -- Present}{M.A. in Mathematics, expected Spring 2027}
\entry{University of Connecticut}{Sept. 2020 -- Dec. 2024}{B.A. in Mathematics}
\entry{East China Normal University}{Sept. 2020 -- Dec. 2021}{Study abroad program during the COVID-19 pandemic}

\section*{Teaching Experience}
\entry{Saint Louis University}{Aug. 2026 -- Present}{Graduate Teaching Assistant; Pre-Calculus Instructor (MATH 1400, Fall 2026)}
\begin{itemize}
\item Teach MATH 1400 Pre-Calculus, using worked examples, individual problem-solving, and small-group discussion.
\item Design lecture notes, worksheets, quizzes, examinations, and solution keys; align weekly assessments with classroom practice.
\item Use Canvas and MyMathLab to organize course materials and assignments; maintain a public teaching portfolio on GitHub.
\end{itemize}
\textbf{Precalculus teaching and course-material preparation}\hfill Spring 2026\par
Prepared quizzes, review materials, examinations, and corresponding solutions.

\section*{Selected Academic Projects}
\textbf{Multivariable Taylor's Theorem}\hfill 2026\par
Developed an exposition and presentation for MATH 5023 on directional derivatives and the multivariable Taylor theorem. Explained the proof by reduction to one variable and applications to error estimates and critical points.

\textbf{Calculating Normal Probabilities without Tables}\hfill Fall 2025\par
Presented a graduate project for MATH 5021 at Saint Louis University on Bagby's 1995 approximation method, using double integrals and polar coordinates; prepared slides and a poster.

\section*{Research Interests}
Artificial intelligence and machine learning, with a current focus on visual representations and image retrieval.

\section*{Technical Skills}
\textbf{Programming and ML:} Python, NumPy, PyTorch, Hugging Face Transformers.\par
\textbf{Methods:} Pretrained vision-model inference, image feature extraction, retrieval evaluation, and numerical linear algebra (SVD and similarity measures).\par
\textbf{Tools:} LaTeX, GitHub, Canvas, and MyMathLab.
\vfill
{\small\color{gray}Updated September 2026}
\end{document}
